% ============================================================================
% Fichier  : partie4/P04_C14_changements_configuration.tex
% Titre    : Changements, Mises en Production et Configuration
% Auteur   : MINKA MI NGUIDJOÏ Thierry Emmanuel
% Version  : 1.0 -- octobre 2026
% ============================================================================
% Sources : ITIL 4 (sources publiques, notions reformulées) ; COBIT 2019,
% fiches BAI06, BAI07 et BAI10 (dossier Normes). Laboratoire associé :
% GCI-LAB-07 (évaluation d'un changement, analyse d'impact, décision).
% ============================================================================

\chapter{Changements, Mises en Production et Configuration}
\label{chap:itil-changements}
\chaptermark{Changements et configuration}

\epigraph{%
    « Toutes les grandes pannes de ce cours ont commencé par un
    changement. »%
}{Leçon des grandes défaillances}

\niveauChapitre{Intermédiaire}{%
    Maîtrise du Chapitre~\ref{chap:itil-incidents} (incidents et
    problèmes) et du Chapitre~\ref{chap:defaillances} (grandes
    défaillances).%
}

\begin{boxobjectifs}
\begin{itemize}
    \item Distinguer changement standard, normal et urgent.
    \item Évaluer le risque d'un changement et organiser son autorisation.
    \item Planifier une mise en production et son retour arrière.
    \item Exploiter une CMDB pour une analyse d'impact.
    \item Relier ces pratiques à l'objectif COBIT BAI06.
\end{itemize}
\end{boxobjectifs}

\bigskip

% ============================================================================
\section*{Question directrice : comment changer sans casser~?}
% ============================================================================

Knight Capital, TSB, CrowdStrike, la BANQUE SANAGA~: le
Chapitre~\ref{chap:defaillances} a montré que chaque grande panne avait
commencé par un changement. Pourtant, une organisation qui ne change
rien meurt aussi sûrement~: le correctif que KRIBILOG n'a jamais installé
(Chapitre~\ref{chap:itil-incidents}) est la cause de 40~\% de ses tickets.
Et la BANQUE SANAGA doit relancer, en 2027, la bascule de son système
central.

Le dilemme est celui de COBIT, qui rattache la gestion des changements à
l'\emph{agilité} (Chapitre~\ref{chap:cobit-objectifs})~: il faut changer
\emph{vite} et changer \emph{bien}. Ce chapitre présente les trois
pratiques qui y répondent~: l'habilitation des changements
(sections~\ref{sec:c14-types} et~\ref{sec:c14-risque}), la gestion des
mises en production et des déploiements (section~\ref{sec:c14-mep}), et la
gestion de la configuration (section~\ref{sec:c14-config}). Le laboratoire
GCI-LAB-07 fait évaluer et décider un changement de la banque.

\separateur

% ============================================================================
\section{Les Changements et leurs Types}
\label{sec:c14-types}
% ============================================================================

\subsection{Qu'est-ce qu'un changement~?}

\begin{boxdef}
\termecle{Changement}\index{changement} (\textit{change})~: ajout,
modification ou retrait de tout élément susceptible d'avoir un effet,
direct ou indirect, sur un service.

\medskip
\termecle{Habilitation des changements}\index{habilitation des
changements} (\textit{change enablement})~: pratique qui vise à maximiser
le nombre de changements réussis, en veillant à ce que les risques soient
correctement évalués, à ce que les changements soient autorisés et à ce
que leur calendrier soit maîtrisé.
\end{boxdef}
{\small\itshape Définitions reformulées d'après ITIL~4~\cite{axelos2019itil4}.
La pratique s'appelait \emph{change control} dans la première édition
d'ITIL~4.}

Le mot \emph{habilitation} n'est pas anodin. ITIL~4 ne cherche pas à
freiner les changements, mais à en permettre le plus grand nombre possible
\emph{avec succès}. Une organisation où chaque changement attend trois
semaines l'avis d'un comité contourne tôt ou tard la procédure~: c'est
l'erreur symétrique de celle de Knight Capital.

\subsection{Trois types de changements}

\begin{tableaucours}{L{2.4cm}L{4.8cm}L{5.0cm}}{%
    \tableauHeader{Type} &
    \tableauHeader{Caractéristiques} &
    \tableauHeader{Exemple}
}
Standard & Risque faible, bien connu, procédure documentée~; autorisé
d'avance & Créer un compte utilisateur, ajouter de la mémoire à un serveur
virtuel \\
Normal & Doit être évalué et autorisé selon son risque, par l'autorité
compétente & Installer le correctif de l'éditeur chez KRIBILOG~; basculer
le système central de la BANQUE SANAGA \\
Urgent & Doit être mis en œuvre au plus vite, par exemple pour résoudre un
incident ou combler une faille~; procédure accélérée & Corriger une faille
de sécurité activement exploitée \\
\end{tableaucours}
\vspace{-0.8em}
\begin{center}{\small\itshape Tableau~14.1. Les trois types de changements}\end{center}

\begin{boxalert}
\textbf{Violation typique.} Pour échapper au comité des changements, une
équipe qualifie systématiquement ses changements d'« urgents ». À la BANQUE
SANAGA, 35~\% des changements de 2025 étaient enregistrés comme urgents~;
aucun n'avait fait l'objet de la revue a posteriori prévue.

\smallskip
\textbf{Règle~:} un changement urgent n'est pas un changement non
contrôlé. Il est autorisé par une autorité désignée à l'avance (souvent
restreinte), documenté, et revu après coup (BAI06.02). Un taux de
changements urgents élevé est un indicateur d'alerte.
\end{boxalert}

\subsection{L'autorité de changement}

ITIL~4 parle d'\termecle{autorité de changement}\index{autorité de
changement} plutôt que de comité~: la personne ou le groupe qui autorise
un changement. Elle n'est pas unique~: l'autorisation est
\emph{déléguée} selon le risque. Un changement standard est autorisé
d'avance~; un changement normal à faible risque peut l'être par le
responsable de l'équipe~; un changement à risque élevé exige un
\termecle{comité des changements}\index{comité des changements}
(\textit{change advisory board}, \sigle{CAB}), voire la direction.

Cette délégation est une question de \textbf{droits de décision}
(Chapitre~\ref{chap:alignement})~: qui peut décider de quel changement.
Elle doit être fixée par écrit et approuvée par la gouvernance.

% ============================================================================
\section{Évaluer et Autoriser}
\label{sec:c14-risque}
% ============================================================================

\subsection{Une grille d'évaluation du risque}

Chaque organisation définit sa grille. Celle que retient la BANQUE SANAGA
après la panne note six critères de~1 (favorable) à~3 (défavorable)~; le
total, de 6 à~18, détermine l'autorité compétente.

\begin{tableaucours}{L{3.6cm}L{2.8cm}L{2.8cm}L{2.8cm}}{%
    \tableauHeader{Critère} &
    \tableauHeader{1} &
    \tableauHeader{2} &
    \tableauHeader{3}
}
Criticité du service & Soutien & Important & Critique (paiements,
système central) \\
Utilisateurs touchés & Moins de 50 & 50 à 500 & Plus de 500 ou clients \\
Complexité & Un composant & Plusieurs composants & Plusieurs systèmes,
données converties \\
Réversibilité & Retour arrière testé & Retour arrière prévu non testé &
Retour arrière impossible \\
Couverture des tests & Plus de 95~\% & 80 à 95~\% & Moins de 80~\% \\
Expérience & Changement déjà réalisé & Changement analogue & Première fois \\
\end{tableaucours}
\vspace{-0.8em}
\begin{center}{\small\itshape Tableau~14.2. Grille d'évaluation du risque
d'un changement (convention de la BANQUE SANAGA)}\end{center}

\begin{tableaucours}{L{2.4cm}L{2.6cm}L{7.2cm}}{%
    \tableauHeader{Score} &
    \tableauHeader{Risque} &
    \tableauHeader{Autorité de changement}
}
6 à 9 & Faible & Responsable de l'équipe concernée \\
10 à 13 & Moyen & Responsable des changements, après avis du RSSI \\
14 à 18 & Élevé & Comité des changements~; et, si la réversibilité vaut~3,
décision formelle du comité de pilotage avec acceptation écrite du risque \\
\end{tableaucours}
\vspace{-0.8em}
\begin{center}{\small\itshape Tableau~14.3. Score de risque et autorité
compétente}\end{center}

\begin{boxscenario}{Opération SANAGA. La bascule notée}
\textbf{Notation a posteriori.} Criticité~3, utilisateurs~3, complexité~3,
réversibilité~3 (conversion irréversible), tests~3 (63~\%), expérience~3
(première migration). Score~: \textbf{18 sur 18}.

\medskip
\textbf{Analyse} (confiance élevée). Avec la grille, la bascule aurait
relevé du comité des changements \emph{et} d'une décision formelle du
comité de pilotage, avec acceptation écrite du risque. Aucune des deux
instances ne s'est prononcée. La grille n'aurait pas empêché la décision
de basculer, mais elle aurait obligé quelqu'un à la prendre, par écrit,
en connaissance de cause.
\end{boxscenario}

\subsection{Le calendrier des changements}

Le \termecle{calendrier des changements}\index{calendrier des changements}
(\textit{change schedule}) planifie les changements autorisés et en
informe les parties concernées. Il permet d'éviter les collisions (deux
changements sur le même système le même soir) et d'appliquer des
\emph{périodes de gel} pendant les moments sensibles~: la clôture des
comptes pour une banque, la paie des fonctionnaires pour l'ANPP, l'arrivée
d'un grand navire pour KRIBILOG.

\begin{boxobjet}{BAI06 · Gérer les changements informatiques}
\textbf{Finalité~:} permettre la livraison rapide et fiable des
changements, en limitant le risque d'affecter la stabilité ou l'intégrité
de l'environnement modifié.

\textbf{Objectifs liés~:} EG01 (portefeuille compétitif)~; AG06 (agilité).

\textbf{Pratiques~:} BAI06.01 évaluer, prioriser et autoriser les demandes
de changement~; BAI06.02 gérer les changements d'urgence~; BAI06.03 suivre
l'état des changements et en rendre compte~; BAI06.04 clore et documenter
les changements.

\textbf{RACI (extrait, BANQUE SANAGA)~:} BAI06.01 pour un risque élevé~:
\emph{A} au président du comité des changements, \emph{R} au responsable
des changements, \emph{C} au RSSI, aux métiers et à l'exploitation,
\emph{I} au centre de services.

\textbf{Indicateurs~:} taux de changements en échec~; taux de changements
urgents~; nombre de changements non autorisés détectés~; part des
changements à risque élevé dotés d'un retour arrière testé.
\end{boxobjet}
{\small\itshape Finalité, objectifs liés et pratiques résumés et traduits
par nos soins à partir de COBIT~2019~\cite{isaca2018cobitobjectifs}.}

\begin{boxmetier}{Responsable des services IT}
\textbf{La charte du comité des changements}
\begin{enumerate}
    \item \textbf{Mandat}~: autoriser les changements à risque élevé,
          arbitrer les conflits de calendrier, examiner les changements en
          échec.
    \item \textbf{Composition}~: responsable des changements (président),
          exploitation, sécurité, représentants des métiers concernés,
          fournisseur si besoin.
    \item \textbf{Fréquence}~: hebdomadaire, et sur convocation pour un
          changement urgent à risque élevé.
    \item \textbf{Pièces exigées}~: évaluation du risque, plan de mise en
          production, plan de retour arrière et preuve de son test,
          résultats de recette.
    \item \textbf{Trace}~: décision motivée, inscrite au registre des
          changements avant toute mise en production.
\end{enumerate}
\end{boxmetier}

% ============================================================================
\section{Mises en Production et Déploiements}
\label{sec:c14-mep}
% ============================================================================

\subsection{Deux pratiques distinctes}

ITIL~4 distingue la \termecle{gestion des mises en production}\index{mise en
production} (\textit{release management}), qui rend disponibles des
services ou des fonctions nouveaux ou modifiés, et la \termecle{gestion
des déploiements}\index{déploiement} (\textit{deployment management}),
qui transfère les composants vers les environnements cibles. On peut
déployer un logiciel sans le rendre disponible aux utilisateurs, puis
l'activer plus tard~: c'est le principe des fonctions activables à la
demande.

\subsection{Stratégies de mise en production}

\begin{tableaucours}{L{3.0cm}L{4.8cm}L{4.4cm}}{%
    \tableauHeader{Stratégie} &
    \tableauHeader{Principe} &
    \tableauHeader{Risque}
}
Bascule unique (\textit{big bang}) & Tout le monde passe au nouveau système
au même moment & Élevé~: tout échec touche tout le monde \\
Par étapes & Agence par agence, région par région & Modéré~: l'échec est
circonscrit \\
Pilote & Un site ou un groupe d'utilisateurs d'abord & Faible pour le
pilote, mais long \\
Fonctionnement en parallèle & Ancien et nouveau systèmes tournent
ensemble pendant une période & Faible, mais coûteux \\
Déploiement progressif (\textit{canary}) & Une petite part du trafic d'abord,
puis le reste & Faible~; suppose une architecture adaptée \\
\end{tableaucours}
\vspace{-0.8em}
\begin{center}{\small\itshape Tableau~14.4. Stratégies de mise en
production}\end{center}

La BANQUE SANAGA avait choisi une bascule unique, avec une conversion des
données irréversible, pour ses 38~agences. CrowdStrike avait déployé sa
mise à jour en une fois sur tout son parc (Chapitre~\ref{chap:defaillances}).
Dans les deux cas, la stratégie a transformé une erreur en catastrophe.

\subsection{Le retour arrière}

\begin{boxdef}
\termecle{Plan de retour arrière}\index{retour arrière} (\textit{back-out
plan, rollback})~: ensemble des actions qui permettent de revenir à l'état
antérieur au changement si celui-ci échoue, avec les critères qui
déclenchent ce retour et le délai dans lequel il doit être décidé.
\end{boxdef}

Un retour arrière n'existe que s'il a été \emph{testé}. Quand il est
impossible (conversion irréversible des données, mise à jour d'un
schéma de base de données), le risque doit être traité autrement~:
fonctionnement en parallèle, sauvegarde complète vérifiée, ou acceptation
formelle par l'instance compétente.

\subsection{Les pratiques de COBIT}

L'objectif BAI07 (gérer l'acceptation et la transition des changements)
détaille ce qu'il faut maîtriser~: établir un plan de mise en œuvre
(BAI07.01), planifier la conversion des processus, des systèmes et des
données (BAI07.02), planifier et réaliser les tests d'acceptation
(BAI07.03 à BAI07.05), passer en production et gérer les versions
(BAI07.06), assurer un soutien renforcé au démarrage (BAI07.07) et
conduire une revue après la mise en œuvre (BAI07.08)~\cite{isaca2018cobitobjectifs}.
La pratique BAI07.02, la conversion des données, est exactement celle qui
a fait défaut à la BANQUE SANAGA.

\begin{boxprocedure}
\textbf{Les critères « go / no-go ».} Avant toute mise en production à
risque élevé, la décision de passer est prise sur une liste de critères
fixée à l'avance, et consignée~:
\begin{itemize}
    \item recette réalisée au taux convenu, anomalies bloquantes closes~;
    \item retour arrière testé, avec son délai de mise en œuvre~;
    \item données converties contrôlées sur un échantillon, avec un
          rapprochement des soldes~;
    \item équipes de soutien et fournisseur mobilisés pour la période de
          démarrage~;
    \item communication aux utilisateurs prête~;
    \item heure limite au-delà de laquelle le retour arrière est
          automatiquement déclenché.
\end{itemize}
Un seul critère non rempli suffit à décider « no-go »~; le passer outre
exige une acceptation écrite du risque.
\end{boxprocedure}

\begin{boxscenario}{Opération SANAGA. La nouvelle bascule}
\textbf{Plan de 2027.} Bascule par étapes~: une agence pilote à Douala, puis
quatre vagues de huit à dix agences, à un mois d'intervalle. Fonctionnement
en parallèle de l'ancien système pendant deux semaines pour chaque vague,
avec rapprochement quotidien des soldes. Gel des changements pendant la
clôture de fin d'année.

\medskip
\textbf{Évaluation avec la grille.} Criticité~3, utilisateurs~2 (une vague),
complexité~3, réversibilité~1 (fonctionnement en parallèle), tests~1,
expérience~2 (après le pilote). Score~: 12, soit un risque moyen par
vague~; mais le comité de pilotage décide de garder le comité des
changements comme autorité pour chaque vague, compte tenu de l'histoire
de la banque.
\end{boxscenario}

% ============================================================================
\section{La Gestion de la Configuration}
\label{sec:c14-config}
% ============================================================================

\subsection{Éléments de configuration et CMDB}

\begin{boxdef}
\termecle{Élément de configuration}\index{élément de configuration}
(\textit{configuration item}, \sigle{CI})~: tout composant qui doit être
géré pour fournir un service~: un serveur, une application, une base de
données, un contrat, une documentation.

\medskip
\termecle{Base de données de gestion de la configuration}\index{CMDB}
(\textit{configuration management database}, \sigle{CMDB})~: base qui
enregistre les éléments de configuration, leurs attributs et surtout
leurs \emph{relations} (dépend de, est hébergé sur, utilise).
\end{boxdef}

La gestion de la configuration se distingue de la gestion des actifs
(BAI09)~: l'actif s'intéresse à la valeur et au cycle de vie financier
d'un bien (achat, amortissement, mise au rebut)~; la configuration
s'intéresse à son rôle dans le service et à ses relations avec les autres
composants.

\subsection{L'analyse d'impact}

La valeur de la CMDB est de répondre, avant un changement, à la question
« qu'est-ce qui sera touché~? ».

\begin{boxscenario}{Opération KRIBI. Installer le correctif}
\textbf{Extrait de la CMDB pour le CI « Logiciel de suivi des conteneurs,
version~3.2 ».}

\medskip
\begin{tableaucours}{L{4.4cm}L{2.6cm}L{5.0cm}}{%
    \tableauHeader{Élément lié} &
    \tableauHeader{Relation} &
    \tableauHeader{Conséquence pour le changement}
}
Serveur d'application SRV-KRI-02 & Hébergé sur & Redémarrage nécessaire \\
Base de données des mouvements & Utilise & Modification du schéma~:
sauvegarde obligatoire \\
Passerelle des balises GPS & Reçoit de & Test de reconnexion requis \\
Portail clients de l'armateur & Fournit à & Interruption à annoncer au
client \\
Contrat NETPORT SERVICES & Exploité par & Intervention de l'infogérant à
planifier \\
\end{tableaucours}

\medskip
\textbf{Évaluation.} Criticité~3, utilisateurs~3, complexité~2,
réversibilité~2 (modification du schéma de la base), tests~2, expérience~2.
Score~14~: risque élevé, autorité~: comité des changements. Fenêtre
retenue~: dimanche de 2~h à 6~h, hors arrivée de navire.

\medskip
\textbf{Ce que révèle l'exercice.} Sans la CMDB, la modification du schéma
de la base et l'effet sur le portail de l'armateur seraient passés
inaperçus. C'est précisément ce que KRIBILOG ignorait lorsqu'elle a
renoncé, il y a six mois, à installer le correctif.
\end{boxscenario}

\subsection{L'objectif BAI10}

COBIT consacre à la configuration l'objectif BAI10 (gérer la
configuration), dont les pratiques couvrent la définition d'un modèle de
configuration, la tenue d'un référentiel et de ses versions de référence,
la maîtrise des éléments de configuration, les états de configuration, et
la vérification périodique de l'intégrité du référentiel. Ses liens
principaux sont le risque métier (EG02), la continuité (EG06) et la
sécurité (AG07)~\cite{isaca2018cobitobjectifs}~: une CMDB fausse est un
risque, car elle conduit à des analyses d'impact fausses.

\begin{boxalert}
\textbf{Violation typique.} Une organisation achète un outil de CMDB, le
remplit une fois par un inventaire automatique, puis ne le met plus à jour.
Deux ans plus tard, la moitié des relations sont fausses, et plus personne
ne s'y fie.

\smallskip
\textbf{Règle~:} la CMDB se met à jour par le processus de changement~:
tout changement clos met à jour les éléments de configuration concernés
(BAI06.04, BAI10.03), et un contrôle périodique vérifie son intégrité
(BAI10.05).
\end{boxalert}

% ============================================================================
\section{Changements et Gouvernance}
\label{sec:c14-gouvernance}
% ============================================================================

\begin{boxmetier}{Auditeur SI}
\textbf{Détecter les changements non autorisés}
\begin{enumerate}
    \item Extraire les journaux de la production (déploiements, modifications
          de configuration) sur la période.
    \item Les rapprocher du registre des changements~: tout déploiement
          sans changement autorisé correspondant est un écart.
    \item Sur un échantillon de changements autorisés, vérifier que
          l'autorisation est antérieure à la mise en production, que les
          tests et le retour arrière sont documentés.
    \item Vérifier la séparation des tâches~: celui qui développe ne
          déploie pas seul en production.
\end{enumerate}
Ce rapprochement est le test le plus puissant des contrôles généraux sur
les changements~: il part de ce qui s'est réellement passé, et non de ce
qui a été déclaré.
\end{boxmetier}

\begin{boiteRemarque}{DevOps tue-t-il la gestion des changements~?}
Les organisations qui déploient plusieurs fois par jour ne peuvent pas
réunir un comité avant chaque mise en production. Elles ne renoncent pas
pour autant au contrôle~: elles le \emph{déplacent} dans la chaîne de
déploiement automatisée (tests obligatoires, revue de code par un pair,
déploiement progressif, retour arrière automatique) et traitent la plupart
de leurs changements comme des changements standard. Le
Chapitre~\ref{chap:agile-cloud} y reviendra. La gouvernance n'exige pas un
comité~; elle exige que le risque soit évalué et que l'autorisation soit
traçable.
\end{boiteRemarque}

% ============================================================================
\separateur
% ============================================================================

\begin{boxresume}
\textbf{Ce qu'il faut retenir de ce chapitre~:}
\begin{itemize}
    \item \textbf{Changer vite et bien}~: l'habilitation des changements
          vise à maximiser le nombre de changements réussis, pas à les
          freiner.
    \item \textbf{Trois types}~: standard (autorisé d'avance), normal
          (évalué et autorisé selon son risque), urgent (accéléré, mais
          contrôlé et revu après coup).
    \item \textbf{L'autorité de changement est déléguée selon le risque}~:
          c'est une question de droits de décision, fixée par la
          gouvernance.
    \item \textbf{La stratégie de mise en production} détermine l'ampleur
          d'un échec~; un retour arrière n'existe que s'il a été testé.
    \item \textbf{Les critères go / no-go} sont fixés à l'avance~; les
          passer outre exige une acceptation écrite du risque.
    \item \textbf{La CMDB} sert l'analyse d'impact~; elle se tient à jour
          par le processus de changement.
    \item \textbf{BAI06, BAI07, BAI10} sont les objectifs COBIT
          correspondants~; le rapprochement entre journaux de production et
          registre des changements en est le meilleur test d'audit.
\end{itemize}
\end{boxresume}

\bigskip

\begin{boxexo}[title={\ding{45}\quad Exercice~14.1. Quel type de changement~? \niveaubase}]
Qualifiez chacun des changements suivants (standard, normal, urgent) et
justifiez.
\begin{enumerate}[(a)]
    \item Ajouter un utilisateur au groupe « consultation » du logiciel de
          suivi.
    \item Corriger, un samedi soir, une faille activement exploitée sur le
          portail de paiement de l'ANPP.
    \item Migrer la messagerie de la BANQUE SANAGA vers un service en
          nuage.
    \item Remplacer un disque défectueux dans une baie redondante.
\end{enumerate}
\end{boxexo}

\begin{boxexo}[title={\ding{45}\quad Exercice~14.2. Évaluer un changement \niveaumoyen}]
L'ANPP veut mettre en production une nouvelle version de son module de
paie, trois jours avant la paie des fonctionnaires. La recette couvre
85~\% des scénarios, le retour arrière est prévu mais non testé, le
changement touche 400~utilisateurs internes et tous les agents de l'État.
\begin{enumerate}[(a)]
    \item Notez le changement avec la grille du Tableau~14.2 et déduisez
          l'autorité compétente.
    \item Quels critères go / no-go ne sont pas remplis~?
    \item Que recommandez-vous, et qui doit le décider~?
\end{enumerate}
\end{boxexo}

\begin{boxexo}[title={\ding{45}\quad Exercice~14.3. Refonder la gestion des changements de la BANQUE SANAGA \niveauexpert}]
Le comité stratégique des SI demande un dispositif de gestion des
changements qui atteigne le niveau de capacité~3 en dix-huit mois
(Chapitre~\ref{chap:cobit-conception}).
\begin{enumerate}[(a)]
    \item Définissez les trois types de changements de la banque, avec des
          exemples, et la liste initiale des changements standard.
    \item Fixez la délégation de l'autorité de changement, et indiquez qui
          doit l'approuver.
    \item Proposez cinq indicateurs pour le comité d'audit, avec leur cible.
    \item Indiquez comment l'auditrice interne vérifiera, dans un an, que le
          dispositif fonctionne.
\end{enumerate}
\end{boxexo}

\bigskip

\begin{boxinfo}
\textbf{Transition.} Les changements modifient les services~; encore
faut-il savoir ce que l'on a promis à leurs utilisateurs. Le
Chapitre~\ref{chap:itil-niveaux} traite des engagements de service, de la
disponibilité, de la continuité et de la gestion des fournisseurs.
\end{boxinfo}
